\documentclass[10pt,a4paper]{article}
\usepackage[margin=2.2cm]{geometry}
\usepackage{fontspec}
\usepackage{kotex}
\setmainfont{Apple SD Gothic Neo}
\setsansfont{Apple SD Gothic Neo}
\usepackage{amsmath,amssymb}
\usepackage{graphicx}
\usepackage{booktabs}
\usepackage{caption}
\usepackage[dvipsnames]{xcolor}
\usepackage[colorlinks=true,linkcolor=MidnightBlue,citecolor=MidnightBlue,
            urlcolor=MidnightBlue]{hyperref}
\usepackage{titlesec}
\titleformat{\section}{\large\bfseries}{\thesection}{0.6em}{}
\titleformat{\subsection}{\normalsize\bfseries}{\thesubsection}{0.5em}{}
\captionsetup{font=small,labelfont=bf}
\setlength{\parskip}{0.35em}
\setlength{\parindent}{0pt}

\usepackage{mdframed}
\newmdenv[linewidth=0.4pt,linecolor=gray!60,backgroundcolor=gray!5,
          innertopmargin=6pt,innerbottommargin=6pt,skipabove=8pt,skipbelow=8pt]{claimbox}
\newcommand{\claim}[2]{\begin{claimbox}\textbf{주장.} #1\par\smallskip
  \textcolor{gray!70}{\textbf{반증 조건.} #2}\end{claimbox}}

\title{\bfseries 구십팔 퍼센트를 닫고 문턱 앞에서\\[2pt]
\large 지는 곳은 보정이 없는 셋뿐이다}
\author{월드모델 실험실 · 노트 $801$}
\date{2026-08-07}
\begin{document}\maketitle

\section{한 줄}

노트 $800$ 이 연 문(안 옮김이 밖 옮김을 가리킨다 · $r\,0.762$)으로 \textbf{보정을
만들었다}: 유보 예보를 도메인별 안 옮김만큼 미리 밀고 등백분위로 되돌린다 ---
\textbf{전부 학습에서 나오므로 한 건씩 가능하다.}

\begin{center}
\begin{tabular}{lrr}
\toprule
& 자릿수 오차 & 중앙절대오차 \\
\midrule
무보정 등백분위 & $0.4367$ & $0.9615$ \\
\textbf{보정 등백분위} & $\mathbf{0.1713}$ & $\mathbf{0.4833}$ \\
기후값 & $0.1656$ & $0.4855$ \\
유보백분위(전이적 참고) & $0.1880$ & $0.5081$ \\
\bottomrule
\end{tabular}
\end{center}

\textbf{무보정 격차의 $98\%$ 가 닫혔고, 전이적 상한($0.1880$)보다도 낫고,
보조자에서는 기후값도 이긴다.} 주자의 판 가중에서만 기후값에 $0.0057$ 진다 ---
그리고 그 패배는 \textbf{보정이 없는 세 도메인이 통째로 만든 것이다.}

\begin{figure}[t]\centering
\includegraphics[width=0.99\textwidth]{figs/corr.pdf}
\caption{왼쪽: 네 팔. 가운데: \textbf{B$-$C 도메인별} --- 지는 곳(빨강)은 전부
무보정. 오른쪽: 가용성이 가른 두 무리.}
\end{figure}

\section{지는 곳은 보정이 없는 셋뿐이다}

노트 $800$ 이 안 옮김을 못 재는 도메인 셋(도서 행 부족 · 시장팝업·팝업 훈련 밖)은
사전등록대로 무보정으로 뒀다. 그 셋에서 B$(=$A$)$ $0.4079$ 대 C $0.1498$ ---
여기가 판 가중을 뒤집었다.

\begin{center}
\textbf{보정이 덮는 $8$개 도메인(행 $3{,}015$)만 모으면}\\
B $\mathbf{0.1434}$ 대 C $0.1675$ --- \textbf{B 가 $0.0241$ 이긴다.}\\
게임($-0.128$) · 애니($-0.052$) · 웹툰($-0.019$)은 씨앗 $3$SD 밖 승.
\end{center}

\claim{\textbf{예보 기반 절대값이 기후값을 이긴 첫 무리다} --- 이 세션 내내
(논문 $159 \cdot 162$) 절대값은 어느 자로도 기후값을 못 넘었다. 그리고 이 $8$개
묶음은 \textbf{결과가 아니라 가용성}(노트 $800$ 의 행수 규칙)으로 정해졌다 ---
결과를 보고 고른 무리가 아니다.}
{\textbf{그래도 승격하지 않는다.} $8$개 묶음 채점은 사전등록에 없던 절단이고,
씨앗 $0{\sim}3$ 으로 고른 것을 같은 씨앗으로 확인할 수는 없다. \textbf{노트
$802$ 를 등록했다} --- $8$개 목록을 못박고, 새 씨앗($4{\sim}7$), \textbf{실용
동점띠 $0.01$} 로 다시 잰다. (가)면 능력 카드에 \textbf{처음으로 한 건씩
절대값의 조건부 허용}이 올라간다.}

\section{규칙 결함 넷째 --- 이번엔 갈래끼리 충돌했다}

사전등록 (라)는 \emph{``B 가 A 보다 나쁘거나 C 보다 $3$배 밖으로 나쁨 $\to$ 보정
금지''}였다. \textbf{B$-$C $= +0.0057$ 이 잡음의 $4.1$배라 (라)가 문자 그대로
발동한다} --- 그런데 B 는 A 를 $0.2654$ 개선하고 D 도 이긴다. \emph{보정 금지}라는
딱지가 자료와 정반대다.

\claim{\textbf{동점띠가 붕괴했다.} 기후값은 결정론적(씨앗 SD $0$)이라 $3\times$
씨앗SD 띠가 $0.0042$ 로 줄었고, \textbf{$0.57$\%p 차이가 ``$4.1$배 밖''}이 됐다.
논문 $159$ 에서는 두 팔 다 확률적이라 띠가 $0.0275$ 였다 --- 같은 규칙이 기준선의
성질에 따라 $6.5$배 다른 엄격함이 된다. \textbf{규약 $46$}: \emph{결정론적
기준선과 견줄 때 동점띠는 실용 폭으로 미리 정한다. 그리고 갈래를 배타로 짤 때
우선순위를 정한다} --- (라)가 (다)를 가려 도메인 셋의 $3$SD 밖 승리가 판정문에서
사라질 뻔했다.}
{규약 $45$(네 방향)를 지키고도 구멍이 났다 --- 구멍의 층이 올라가고 있다
(갈래 부재 $\to$ 갈래 상호작용). \textbf{다음 층은 아마 ``자 선택''일 것이다}
(주자를 자릿수로 못박은 것 자체가 이번에 판정을 갈랐다 --- 보조자로는 이겼다).
그것을 미리 적어 둔다.}

\section{예측 채점 --- 좋은 쪽으로 틀렸다}

① B $0.19{\sim}0.28$ 로 예측 $\to$ \textbf{$0.1713$ --- 더 좋았다.} ⑤ $|$B$-$D$|
< 0.03$ $\to$ 맞음($0.0167$). ④ \emph{웹툰·펀딩은 덜 닫힐 것}(안 옮김이 하한이라)
$\to$ \textbf{틀렸다} --- 웹툰 $0.042$ · 펀딩 $0.108$ 로 잘 닫혔다.

\claim{\textbf{보정이 전이적 상쇄보다 나은 것이 단서다.} D 는 유보 예보 분포를
그대로 믿는데 B 가 그것을 이겼다 --- \textbf{안 옮김 보정이 드리프트만이 아니라
유보 분포의 잡음 일부까지 걸러낸 것}으로 읽힌다. 그리고 등백분위가 순위 기반이라
보정 크기의 오차(웹툰·펀딩의 하한 문제)에 둔감했다.}
{한 번의 관측이다. 노트 $802$ 의 새 씨앗에서 B$-$D 부호가 뒤집히면 이 해석은
버린다.}

\end{document}
